\section{Sign of the Times}

\begin{quotex}
We must be ready for an immense event in the divine order which we are travelling towards with an accelerated speed that must astound all those who watch. Awesome oracles have already pronounced that the time is now. \flright{\textsc{Joseph de Maistre}}

\end{quotex}
\textbf{Rene Guenon} closed the King of the World with the above quote. The point is not that de Maistre is a fortune teller, but rather he could read the signs of the times. In Chapter IX of the Reign of Quantity, Guenon explains how to become ready, when he writes:

\begin{quotex}
Those who have not effectively attained to such a state must at least, as far as their capabilities permit, use every endeavor to reach it. 

\end{quotex}
The state he is referring to is the supra-individual state, as he explains:

\begin{quotex}
The being who has attained a supra-individual state is, by that fact alone, released from all the limiting conditions of individuality, that is to say, he is beyond the determinations of ``name and form" which constitute the essence and the substance of his individuality as such. 

\end{quotex}
Clearly, there is some frustration possible, since there is no instruction on how to reach that state, nor a definitive description of it, yet we are supposed to be single-pointedly dedicated to it. So men begin to wander aimlessly, looking here or there for an answer, defending their favorite theories, attacking those they dislike. They seek in vain for an initiation, because of this situation that Guenon describes:

\begin{quotex}
Organizations provide an effective initiation only where an authentic traditional doctrine still exists, offering but a shadow when the spirit of the doctrine no longer vivifies the symbols that are merely its external representation; this happens when the conscious connection with the supreme spiritual centre of the world is finally broken. 

\end{quotex}
Can we still be satisfied with that? Has not Guenon taught us traditional doctrines? The centre is everywhere, so that conscious connection can be instantly re-established. After all, Guenon always takes pains to point out that he is describing states of being, not places. So Guenon offers this encouragement:

\begin{quotex}
It is important to realize that we should be talking of something that is hidden rather than truly lost, because it is not lost to everybody but still possessed in its entirety, albeit by a very few. This gives rise to the possibility for other to rediscover it, provided they search in the proper manner. 

\end{quotex}
In the Ancient City, it was often claimed that the founder had hidden the city's constitution in a tree trunk, or some such place. Finding a secret text is not what is meant, but rather the secret is unveiled in consciousness provided one has the right intention and is endeavoring.

Guenon hints that effective initiatory knowledge can be found in Tibet. He points to Swedenborg as authority in this regard as well as the `Mount of Prophets' described by Anne-Catherine Emmerich. He alludes also to Agarttha as described in Saint-Yves d'Alveydre's Mission of India.

Meanwhile, there is no excuse to not make endeavors. Learn the doctrines, practice the many suggestions in the spiritual literature we have recommended (or your own favorite), and learn to recognize the different states of being in you own consciousness.



\flrightit{Posted on 2013-01-09 by Cologero }

\begin{center}* * *\end{center}

\begin{footnotesize}\begin{sffamily}



\texttt{Mihai on 2013-01-10 at 05:06 said: }

Swendenborg is the ``father" of spiritism as far as I know about him. Where did Guenon ever refer to him in his works ?

Or maybe we are referring to different persons? Did you mean Emmanuel Swendenborg ?


\hfill

\texttt{Cologero on 2013-01-10 at 07:33 said: }

Oddly enough, Guenon refers to Swedenborg often. This specific reference was from the King of the World. In the Spiritist Fallacy, Guenon does mention Swedenborg in relation to the spiritists. However, as he points out, Swedenborg is to be understood symbolically, not literally. For example he writes:

\begin{quotex}
When seers are orthodox mystics, their natural tendencies to stray are in some manner held in check and reduced to a minimum; almost everywhere else they have free rein and the result is often a nearly inextricable confusion. The most unquestionable and most celebrated among them, Swedenborg for example, are far from exempt from this fault, and one cannot take too many precautions if one wishes to extract what is of genuine interest in their works. Better to go to purer sources … 

\end{quotex}
Nevertheless, Guenon does not restrict himself totally to those purer sources. On a personal level, we can presume that Guenon became familiar with Swedenborg from his early involvement with Martinism. Apparently, he felt Swedenborg, in this instance, offered something of interest. Similar things could be said about another seer, Emmerich, who, by the way, served as the inspiration for some scenes in Gibson's the Passion of the Christ. I suppose the point here is to not be too dogmatically close-minded.


\hfill

\texttt{Iulianus on 2013-01-10 at 08:38 said: }

I think that Yves' book Mission de l'Inde was translated and published in English under the title The Kingdom of Agarttha: A Journey into the Hollow Earth…


\hfill

\texttt{Cologero on 2013-01-10 at 12:53 said: }

Thanks for pointing this out, Iulianus. Has anyone read it?


\hfill

\texttt{Mihai on 2013-01-10 at 15:11 said: }

I see.

Oddly, I do not recall any mention in the King of the World regarding him. 

I understand the case, though.

My take in these cases: if you find a piece of gold, use it, without concerning yourself with where it came from. But don't go burying yourself neck-deep in trash because you know that there has to be a piece of gold there, especially if you've got an entire goldmine near you.


\hfill

\texttt{Pickman on 2013-01-10 at 19:15 said: }

``Meanwhile, there is no excuse to not make endeavors. Learn the doctrines, practice the many suggestions in the spiritual literature we have recommended (or you own favorite), and learn to recognize the different states of being in you own consciousness."

One can do these exercises on their own, but without a guiding hand – be that of a master and a structured community of disciplined adherents (often recommended even by Steiners' standards) dangers can and will persist. A student may be off by a millimeter and yet not ever be corrected on the trajectory. This is why institutions existed in the past to buffer such spiritual endeavors for those who had the vocation to commit to them. Guénon, Eckhart, tomberg and other esoteric maistres may have as individuals labored endearingly for years but they still needed to graft themselves to a Tradition for anchorage. Unfortunately we have been left smashed up upon the ruins, to raise ourselves up without a wall to lean on (of the hypertrophied, post-post-modern spiritually vacant, materialist West – notwithstanding persistent witch cults of a feminine neo-spiritism).

``Guenon hints that the effective initiatory knowledge can be found in Tibet. He points to Swedenborg as authority in this regard as well as the `Mount of Prophets' of Anne-Catherine Emmerich. (If anyone has located either of these specific references, please share them.) He alludes also to Agarttha as described in Saint-Yves d'Alveydre's Mission of India. I will be making his description of Agarttha available in English for the first time in the near future."

Was this hinting done in Guénon's theosophy days in his exploration of occultist circles and the claims of Bo yin ra? If at a later time he could not get to India, then surely Tibet was out of the question. Then again there are hints towards initiatic centers even in the West still, yet we are told they have refused to act for reasons known to themselves.Telling really. Which is true even of those few who are receptive to a rightist doctrine. 

Who are the few and who have remembered so that we can have a living example to aim for? That is what is missing today (among other things).


\hfill

\texttt{Senko on 2013-01-10 at 20:42 said: }

I know this can only satisfy Christians (Catholics and Orthodox alike), but Prayer of the Heart and Lectio Divina along with a good Confessor and a Priest that can act as a spiritual guide/mentor, should be enough to give the first steps into spiritual realization. What I find truly difficult in this turbulent times is not finding the appropiate way, but keeping up with it in these freaking crazy times


\hfill

\texttt{Cologero on 2013-01-10 at 23:16 said: }

Guenon hints at this often enough: the goal is not really to ``find" a spiritual center some place not least of all because it ``is to all intents and purposes inaccessible and invisible to all except those possessing the necessary qualifications for entry."

Also, ``both geographical and historical facts possess a symbolic validity that in now way detracts from their being facts."

In other words, the spiritual center that people are seeking (externally) can be right here, right now … it would still be unrecognizable without the necessary qualifications. The few have been born at this era, some of whom will try to remember why. They are the qualified, they are the living examples. There is no magic group to find and you can't wait for George to do it for you. If the world is in ruins … what would a man do?


\hfill

\texttt{Jason-Adam on 2013-01-11 at 06:38 said: }

The Ladder of Divine Ascent is another useful text of the Christian path….


\hfill

\texttt{Jason-Adam on 2013-01-11 at 06:44 said: }

It's funny, I just finished reading Saint-Yves d'Alveydre's book Mission des Souverains and sadly was disappointed by his pro-Masonic and anti-Catholic stance, though I do believe his idea of synarchy once adopted to Catholic means is a good system. I'm interested to read Cologero's translations of him soon…


\hfill

\texttt{Pickman on 2013-01-11 at 14:27 said: }

Reference was made to Tibet on account of them continuing a living pre-modern solar initiation that would be acceptable to those ``possessing the necessary qualifications for entry" and not just hollow rites based on empty form and ritual of decadent Western institutions. The reference was not made on ``finding" or ``seeking" something external or material even beyond the symbolic significance of their locations. An individual of Indo-European heritage already is who he is, but he must activate his center and know himself. This is where it gets confusing and there are many pretenders to the throne. Try joining the Jesuits but it will not be enough for some, but at least a Ur group would have been useful for research purposes. Who's George? 

``What I find truly difficult in this turbulent times is not finding the appropiate way, but keeping up with it in these freaking crazy times"

The only way out is to run vertically into the mountains or horizontally to the wilderness.


\end{sffamily}\end{footnotesize}
